\section{Whole module sections on an original affine cartesian square}

\begin{lemma}\label{lem:fc-affine-cartesian-module-base-change}
\lean{EnvelopeProperty.AffineCartesianModuleBaseChange.tensorEquiv,EnvelopeProperty.AffineCartesianModuleBaseChange.tensorEquiv_tmul}
\leanok
\uses{lem:fc-affine-module-flat-base-change}
Let $R\to S$ be an original homomorphism of commutative rings and let
$g:Y\to X$, $p:X\to\operatorname{Spec}R$ and
$q:Y\to\operatorname{Spec}S$ form an actual cartesian square over this
homomorphism. Suppose the original $X$ and $Y$ are affine. For every
original quasicoherent module $M$ on $X$, the original coefficient maps
and original pullback give an equivalence
\[
 S\otimes_R\Gamma(X,M)\simeq_S\Gamma(Y,g^*M).
\]
Every original pure tensor $s\otimes m$ maps to the actual coefficient
$q^*(s)$ times the original pullback-adjunction-unit section of $m$.
No flatness or section comparison certificate is required.
\end{lemma}
\begin{proof}
The genuine affine cartesian square makes the entire square of original
coordinate rings a pushout. Transport the two spectrum coordinate
rings through their actual spectrum section isomorphisms. Thus, with
$A=\Gamma(X,\mathcal O_X)$ and $B=\Gamma(Y,\mathcal O_Y)$, the actual
ring $B$ is the pushout of $A$ and $S$ over $R$, with their literal
coefficient maps. The square also proves the required two scalar tower
identities. Tensor cancellation for this actual pushout identifies
$S\otimes_R\Gamma(X,M)$ with $B\otimes_A\Gamma(X,M)$. The native
quasicoherent affine pullback equivalence identifies the latter whole
module with $\Gamma(Y,g^*M)$. Its value on $1\otimes m$ is the original
pullback unit. Both maps are linear over the actual coefficient rings,
so the tensor cancellation formula and this native unit formula prove
the assertion on every $s\otimes m$.
\end{proof}
